\subsection{Normalización + Comparación con el Reverso (palíndromos)}
\label{alg:9-2}

\graybold{Preguntas guía:}

\begin{itemize}
\item ¿Piden verificar si un texto es un palíndromo ignorando mayúsculas, espacios y puntuación?
\item ¿Basta con normalizar el string (filtrar solo los caracteres relevantes, unificar case) y comparar contra su reverso?
\item ¿El caso de un string vacío (tras filtrar) debe considerarse un palíndromo válido?
\end{itemize}

\graybold{Utilidad:} Resolver directamente problemas de verificación de palíndromos con reglas de normalización, filtrando primero y comparando el resultado contra su versión invertida.

\graybold{Complejidad:} \bigO{n}.

\graybold{Fundamento teórico:} Un string es palíndromo si es igual a su reverso. Al normalizar primero (quedarse solo con letras, todas en minúscula), la comparación directa \texttt{filtered == filtered[::-1]} captura exactamente la definición pedida, incluyendo el caso borde de que un string vacío (o sin letras) es trivialmente un palíndromo.

\graybold{Ejercicio aplicado}

Dada una oración, determinar si es un palíndromo ignorando mayúsculas, espacios, y cualquier carácter que no sea una letra del alfabeto inglés.

\graybold{I/O:} Múltiples líneas hasta el centinela "*" (patrón 3), N < 1000 por línea → se recorre línea por línea con \texttt{for line in sys.stdin}, ya que es texto libre (no se puede tokenizar con \texttt{.split()}).

\graybold{Código solución}

\begin{lstlisting}[language=Python]
import sys

def solve():
    out = []
    for raw in sys.stdin:
        line = raw.rstrip('\n').rstrip('\r')
        if line == '*':
            break
        filtered = [c.lower() for c in line if c.isalpha()]
        out.append('Y' if filtered == filtered[::-1] else 'N')
    sys.stdout.write("\n".join(out) + "\n")

solve()
\end{lstlisting}
